\section{O Viés Indutivo na Classificação de Granulação Fina ($H_1$)}
\label{sec:vies_indutivo}

A superioridade estatística das Redes Neurais Convolucionais (CNNs) modernas sobre os modelos baseados em \textit{Transformers} (ViT e Swin), evidenciada na Fase 1, revela uma característica fundamental do domínio de aplicação. A classificação de pluma de algodão configura-se como um problema de Classificação de Granulação Fina, onde a distinção entre classes comerciais (e.g., 31.3 vs. 31.4) não reside na estrutura semântica global do objeto, mas em variações sutis de microtexturas, reflectância e densidade de impurezas.

A dominância de modelos como a ConvNeXt e a DenseNet pode ser explicada pelo viés indutivo de localidade inerente às operações convolucionais. Diferentemente dos \textit{Vision Transformers}, que utilizam o mecanismo de auto-atenção global para correlacionar todos os recortes (\textit{patches}) da imagem simultaneamente, as CNNs operam através de campos receptivos locais que são hierarquicamente agregados. Para a tipagem do algodão, onde o diagnóstico final depende de elementos estocásticos e espacialmente distribuídos, como nós de fibras, fragmentos de cascas e estrias de coloração, a capacidade das convoluções em extrair padrões de alta frequência espacial mostrou-se superior à modelagem de dependências de longo alcance.

Adicionalmente, os resultados sugerem que a ausência de um viés espacial forte nos \textit{Transformers} puros exigiria um volume de dados massivamente superior para que a rede convergisse e aprendesse, de forma autônoma, a focar em texturas de pequena escala. No contexto desta pesquisa, as arquiteturas convolucionais modernas demonstraram maior eficiência de dados, extraindo representações latentes mais ricas e discriminativas a partir do conjunto de dados disponível. Este achado atesta que, para o reconhecimento de texturas orgânicas complexas e padronizadas, a topologia convolucional permanece como o estado da arte em termos de assertividade diagnóstica e economia paramétrica. Essa robustez é explicada pela capacidade das CNNs em atuar como detectores hierárquicos de textura \cite{zeiler2014visualizing}. A literatura de explicabilidade, notadamente por meio de técnicas como Grad-CAM \cite{selvaraju2017gradcam}, demonstra que redes convolucionais profundas priorizam gradientes de borda e oclusões espaciais na formação da decisão. No domínio do algodão, isso garante que o modelo capture as micro-impurezas e os nós de fibras que definem o Tipo comercial, minimizando o risco de aprendizado espúrio baseado em características do fundo da caixa de captura.